% BEGIN VOL07-EXPANSION VII38
\section{Supplementary worked examples}
\label{sec:vii38-pedagogy}

\begin{example}[Graph path versus surface shortcut]\label{ex:vii38-ped-01}
Consider a square triangulated by one diagonal.  Between opposite vertices not joined by that diagonal, an edge-graph path follows two sides with length \(2\), while the continuous surface geodesic crosses the flat square with length \(\sqrt2\).  Exact graph search can therefore be geometrically inexact.
\end{example}

\begin{example}[Unfolding two triangles]\label{ex:vii38-ped-02}
When a shortest path crosses a shared interior edge away from vertices, unfold one triangle into the plane of the other.  The two path segments become one straight segment; any kink would admit a shorter chord.
\end{example}

\begin{example}[Angle defect on a tetrahedron]\label{ex:vii38-ped-03}
At each vertex of a regular tetrahedron, three equilateral face angles meet, so \(\Theta=\pi\) and the discrete Gaussian curvature is \(K_v=2\pi-\pi=\pi\).  The four defects sum to \(4\pi\), matching Gauss--Bonnet for a sphere topology.
\end{example}

\section{Graded supplementary exercises}
The following exercises add a deliberate progression from concrete calculation to structural proof, hypothesis testing, application, and synthesis.

\subsection*{Standard computations and constructions}

\begin{exercise}[Square graph]\label{exr:vii38-ped-01}
On a unit square with only boundary edges available, what is the graph distance between opposite corners?
\end{exercise}

\begin{hint}
Two unit edges are required.
\end{hint}

\begin{solution}
\(2\).
\end{solution}

\begin{exercise}[Square surface]\label{exr:vii38-ped-02}
What is the intrinsic flat-surface distance between those corners?
\end{exercise}

\begin{hint}
Use the Euclidean diagonal.
\end{hint}

\begin{solution}
\(\sqrt2\).
\end{solution}

\begin{exercise}[Upper bound]\label{exr:vii38-ped-03}
Which is larger in general for mesh vertices: edge-graph distance or continuous polyhedral distance?
\end{exercise}

\begin{hint}
Every edge path is a valid surface path.
\end{hint}

\begin{solution}
\(d_M\le d_E\).
\end{solution}

\begin{exercise}[Tetra defect]\label{exr:vii38-ped-04}
Compute the angle defect at a regular tetrahedron vertex.
\end{exercise}

\begin{hint}
Three angles of \(\pi/3\) meet.
\end{hint}

\begin{solution}
\(\pi\).
\end{solution}

\begin{exercise}[Total tetra curvature]\label{exr:vii38-ped-05}
Sum all regular tetrahedron angle defects.
\end{exercise}

\begin{hint}
There are four vertices.
\end{hint}

\begin{solution}
\(4\pi\).
\end{solution}

\subsection*{Proofs}

\begin{exercise}[Unfolding condition]\label{exr:vii38-ped-06}
Why must a locally shortest edge-crossing path be straight after unfolding?
\end{exercise}

\begin{hint}
A kink is longer than its chord.
\end{hint}

\begin{solution}
If the unfolded segments meet with nonzero turning angle, replacing them by the straight segment shortens the path while staying in the same faces.
\end{solution}

\begin{exercise}[Intrinsic invariance]\label{exr:vii38-ped-07}
Why does bending an intrinsically isometric triangle mesh in \(\mathbb R^3\) not change its polyhedral geodesic distances?
\end{exercise}

\begin{hint}
The metric is determined by face edge lengths.
\end{hint}

\begin{solution}
Facewise Euclidean lengths and gluing data are unchanged, so all admissible path lengths are unchanged.
\end{solution}

\begin{exercise}[Positive vertex avoidance]\label{exr:vii38-ped-08}
Why do shortest paths generally not pass through a positive-curvature cone tip?
\end{exercise}

\begin{hint}
There is insufficient total angle to place at least \(\pi\) on both sides of a locally straight path.
\end{hint}

\begin{solution}
Unfolding the cone shows a path through the tip can be shortcut unless the two-sided angle condition is met, which positive defects obstruct.
\end{solution}

\begin{exercise}[Distance continuity]\label{exr:vii38-ped-09}
Why can distance remain continuous while its gradient jumps at the cut locus?
\end{exercise}

\begin{hint}
Several minimizing branches can tie in length.
\end{hint}

\begin{solution}
The minimum of smooth branch-distance functions is continuous but nondifferentiable where multiple branches attain the same value.
\end{solution}

\subsection*{Counterexamples and hypothesis tests}

\begin{exercise}[Dijkstra exact surface]\label{exr:vii38-ped-10}
Test the claim on the flat square example.
\end{exercise}

\begin{hint}
Dijkstra is exact only for the graph problem.
\end{hint}

\begin{solution}
False.  It returns \(2\) while the surface distance is \(\sqrt2\).
\end{solution}

\begin{exercise}[Smooth ODE at vertices]\label{exr:vii38-ped-11}
Test whether the smooth Christoffel geodesic equation should be imposed at a cone vertex.
\end{exercise}

\begin{hint}
The piecewise-Euclidean metric is singular there.
\end{hint}

\begin{solution}
No.  Vertex behavior is governed by cone angles and unfolding, not a smooth connection.
\end{solution}

\begin{exercise}[Small scalar error means same path]\label{exr:vii38-ped-12}
Test near a cut locus.
\end{exercise}

\begin{hint}
Different minimizing routes can have nearly identical lengths.
\end{hint}

\begin{solution}
False.  Path geometry can change abruptly while scalar distance error remains tiny.
\end{solution}

\subsection*{Applications and investigations}

\begin{exercise}[Benchmark design]\label{exr:vii38-ped-13}
What two error types should a geodesic benchmark report?
\end{exercise}

\begin{hint}
Separate scalar and geometric path quality.
\end{hint}

\begin{solution}
Distance/length error and a geometric path discrepancy such as sampled Hausdorff distance.
\end{solution}

\begin{exercise}[Source representation]\label{exr:vii38-ped-14}
Why should a point source inside a triangle be stored barycentrically rather than snapped to a vertex?
\end{exercise}

\begin{hint}
Snapping changes the problem before solving it.
\end{hint}

\begin{solution}
Barycentric storage preserves the requested surface point and avoids an uncontrolled source-position error.
\end{solution}

\subsection*{Challenge problems}

\begin{exercise}[Cone criterion]\label{exr:vii38-ped-15}
At an interior vertex with total angle \(3\pi\), classify the curvature sign.
\end{exercise}

\begin{hint}
Compare with \(2\pi\).
\end{hint}

\begin{solution}
It is a saddle vertex with negative angle defect \(-\pi\).
\end{solution}

\begin{exercise}[Refinement challenge]\label{exr:vii38-ped-16}
Why is mesh refinement alone not a guarantee of monotone graph-distance convergence if edge directions remain highly anisotropic?
\end{exercise}

\begin{hint}
Graph paths are restricted to available directions.
\end{hint}

\begin{solution}
Directional bias can persist despite smaller edges; convergence quality depends on both resolution and triangulation geometry.
\end{solution}

% END VOL07-EXPANSION VII38
